\documentclass{article}
\usepackage[T1]{fontenc}
\usepackage{lmodern}
\usepackage{graphicx}
\usepackage{amsmath,amssymb}
\usepackage[a4paper,margin=2cm]{geometry}
\usepackage[absolute,overlay]{textpos}
\setlength{\TPHorizModule}{1mm}
\setlength{\TPVertModule}{1mm}
\usepackage[indonesian]{babel}
\usepackage{hyperref}
\setlength{\emergencystretch}{2em}
% unit: Halaman 20

\begin{document}

% === Halaman 20 (python/native) ===

Contoh 5:  Kita akan mendekripsikan cipherteks dari Contoh 4 dengan 

menggunakan kunci privat   \{2, 3, 6, 13, 27, 52\}. Di sini, n = 31 dan m = 105. 

Nilai 31–1 (mod 105) diperoleh sbb:

  n  n–1  1 (mod m) → 31  n–1  1 (mod 105) → n–1 = (1 + 105k)/31 

    coba k = 0, 1, 2, …, diperoleh n–1 = 61

Cipherteks dari Contoh 4 adalah 174, 280, 333. Hasil dekripsi:

 174  61 mod 105 = 9 = 02 + 13 + 16 + 013 + 027 + 052    → 011000

 280  61 mod 105 = 70 = 12 + 13 + 06 + 113 + 027 + 152 → 110101

 333  61 mod 105 = 48 = 12 + 03 + 16 + 113 + 127 + 052 → 101110

Jadi, plainteks yang dihasilkan kembali adalah: 


 011000110101101110

20

\end{document}
